\documentclass{article}
\usepackage[nonatbib]{neurips_2026}

\usepackage[utf8]{inputenc}
\usepackage[T1]{fontenc}
\usepackage{hyperref}
\usepackage{url}
\usepackage{booktabs}
\usepackage{amsfonts}
\usepackage{amsmath}
\usepackage{nicefrac}
\usepackage{microtype}
\usepackage{xcolor}
\usepackage{graphicx}

\title{Stress-Testing Wanda: A Reproducibility and Robustness Study\\
of Activation-Based LLM Pruning}

\author{%
  Minguhn Kim \\
  Graduate School of Semiconductor Technology \\
  POSTECH \\
  \texttt{minguhn@postech.ac.kr} \\
  \And
  Jihun Park \\
  Graduate School of Convergence Science and Technology \\
  POSTECH \\
  \texttt{jpark42@postech.ac.kr}
}

\begin{document}

\maketitle

\begin{abstract}
We present a systematic stress-test of Wanda~\cite{wanda}, a simple and
effective pruning method for large language models (LLMs) that scores weights
by the product of their magnitude and the $\ell_2$-norm of corresponding
input activations. The original paper claims robustness to the choice of
calibration data, but tests this only by varying the \emph{quantity} of
English C4 samples---never the \emph{domain} or \emph{language}. We argue
that the activation norm $\|X_j\|_2$ is a statistical fingerprint of the
calibration distribution, and that domain-mismatched calibration can make
Wanda perform \emph{worse} than the calibration-free magnitude pruning
baseline on capabilities not represented in the calibration data.
We test this thesis on LLaMA across three experiments:
(1) calibration domain shift at 50\% sparsity using English (C4),
Korean (MC4-ko), code (The Stack), and random-token calibrations;
(2) per-task fragility curves across 50\%, 60\%, and 70\% sparsity;
and (3) the calibration $\times$ sparsity interaction at 70\% sparsity.
We additionally document reproducibility issues encountered when running
the original codebase under modern library versions, and provide an
English-only domain-shift supplement on LLaMA-2-7B to isolate the
cross-lingual component of the effect.
\end{abstract}

% =============================================================
\section{Introduction}
% =============================================================

Model compression is a critical enabler of efficient inference for large-scale
neural networks. Among compression techniques, unstructured pruning offers a
flexible way to reduce model size by zeroing out individual weights.
Wanda~\cite{wanda} (Pruning by \textbf{W}eights \textbf{and} \textbf{A}ctivations)
proposes a remarkably simple pruning criterion: for each weight $w_{ij}$,
its importance score is defined as
\begin{equation}
  S_{ij} = |w_{ij}| \cdot \|\mathbf{x}_j\|_2,
\end{equation}
where $\mathbf{x}_j$ denotes the $j$-th input feature activation collected
over a small calibration set. Weights with the lowest scores are pruned.
This requires no gradient computation and runs in minutes even for
65B-parameter models, making it highly practical.

However, the original evaluation tests calibration-data robustness only by
varying the \emph{quantity} of C4~\cite{c4} samples in a single domain
(English web text). This leaves a sharper question untested: \emph{can
domain-mismatched calibration actively harm Wanda's pruning quality, to the
point of being worse than calibration-free magnitude pruning?}
We focus on this single thesis through three experiments:

\begin{itemize}
  \item \textbf{Calibration domain shift.} At 50\% sparsity, does using
        English, Korean, code, or random-token calibration affect which
        capabilities (English knowledge, math reasoning, Korean language)
        survive pruning?
  \item \textbf{Per-task fragility curves.} As sparsity increases from 50\%
        to 70\%, do different downstream tasks degrade at different
        thresholds, and does perplexity track them?
  \item \textbf{Calibration $\times$ sparsity interaction.} At the
        most aggressive setting, do calibration-domain damage and
        sparsity-cliff damage compound super-additively, suggesting
        both effects target the same low-activation features?
\end{itemize}

Following the Track 2 stress-testing protocol, we first reproduce the
original codebase (and document the reproducibility issues encountered),
then run the three experiments above on LLaMA, with a LLaMA-2-7B
English-only supplement to isolate the cross-lingual component.

% =============================================================
\section{Background}
% =============================================================

\subsection{Wanda Pruning}

Given a pre-trained LLM, Wanda prunes each linear layer independently by
computing importance scores $S_{ij}$ for all weights and zeroing out those
below a target sparsity ratio $s$. The key insight over magnitude-only pruning
is that a weight connected to a high-activation feature should be treated as
more important, even if its magnitude is small. The calibration set is used
solely for a single forward pass to collect activation statistics; no weight
updates are performed.

\subsection{Baselines}

We compare Wanda against:
\begin{itemize}
  \item \textbf{Magnitude pruning}: scores weights by $|w_{ij}|$ only.
  \item \textbf{SparseGPT}~\cite{sparsegpt}: solves a layer-wise
        reconstruction problem; more expensive but generally stronger.
  \item \textbf{Dense baseline}: unpruned model performance.
\end{itemize}

% =============================================================
\section{Reproducibility Issues}
% =============================================================

Before conducting stress-tests, we attempted to reproduce the original results
using the public codebase.\footnote{\url{https://github.com/locuslab/wanda}}
We identified the following compatibility issues with modern library versions
(torch 2.6, transformers 5.5, datasets 3.x):

\begin{enumerate}
  \item \textbf{torch $<$ 2.6 blocked}: A security vulnerability
        (CVE-2025-32434) in \texttt{torch.load} causes the \texttt{transformers}
        library to raise a \texttt{ValueError} when loading model weights with
        torch versions below 2.6. The original codebase was developed with
        torch 2.x, requiring an upgrade.

  \item \textbf{C4 dataset config renamed}: The \texttt{datasets} library
        changed the C4 configuration identifier from \texttt{`allenai--c4'}
        to \texttt{`en'}, causing a \texttt{BuilderConfig} error at runtime.

  \item \textbf{Split size verification}: Even after the config fix,
        loading only a subset of C4 (one shard) triggers a
        \texttt{NonMatchingSplitsSizesError}. The fix requires passing
        \texttt{verification\_mode=`no\_checks'}.

  \item \textbf{OPT architecture mismatch}: Wanda's calibration code
        accesses \texttt{model.model.layers}, which is valid for LLaMA
        but fails for OPT, whose layers reside at
        \texttt{model.model.decoder.layers}.
\end{enumerate}

These issues collectively demonstrate that the original codebase lacks
robustness to software environment changes, raising reproducibility concerns.
All fixes were applied to our fork before proceeding with experiments.

% =============================================================
\section{Experimental Setup}
% =============================================================

\subsection{Compute Environment}

All experiments are conducted on a server equipped with
8$\times$ NVIDIA RTX 6000 Ada Generation GPUs (48\,GB VRAM each),
running CUDA 12.6, PyTorch 2.6.0, and transformers 5.5.3.

\subsection{Models}

\begin{itemize}
  \item \textbf{LLaMA-3-8B}~\cite{llama}: main model. Chosen because it
        scores above random chance on Korean benchmarks (KMMLU, KoBEST),
        which is required to measure Korean-capability loss after pruning.
  \item \textbf{LLaMA-2-7B}~\cite{llama}: supplementary model for the
        English-only domain-shift control (Section~\ref{sec:exp4}).
\end{itemize}

\subsection{Calibration Datasets}

We use 128 samples per source, matching the original Wanda setup:
\begin{itemize}
  \item \textbf{C4}~\cite{c4} (default): English web text.
  \item \textbf{MC4-ko}: Korean subset of the multilingual C4.
  \item \textbf{The Stack}: source-code corpus (a non-natural-language domain).
  \item \textbf{Random tokens}: tokens sampled uniformly from the vocabulary;
        a structureless control.
\end{itemize}

\subsection{Evaluation Metrics}

All metrics are computed with the \texttt{lm-evaluation-harness} framework
where applicable.
\begin{itemize}
  \item \textbf{WikiText-2 perplexity} ($\downarrow$): primary metric in
        the original Wanda paper; English language-modeling fluency.
  \item \textbf{MMLU} ($\uparrow$, 5-shot accuracy): English world-knowledge
        and reasoning.
  \item \textbf{GSM8K} ($\uparrow$, 8-shot accuracy): English mathematical
        reasoning.
  \item \textbf{KoBEST HellaSwag} ($\uparrow$, 0-shot accuracy): Korean
        commonsense reasoning---the capability we expect to be most
        sensitive to English-only calibration.
\end{itemize}

% =============================================================
\section{Results}
% =============================================================

\subsection{Experiment 1: Calibration Domain Shift at 50\% Sparsity}
\label{sec:exp1}

We hold model (LLaMA-3-8B) and sparsity (50\%, unstructured) fixed and vary
the calibration source for Wanda, comparing against two anchor baselines
that do not depend on calibration data the same way.
Table~\ref{tab:exp1-calib-domain} reports six pruning runs.

\begin{table}[h]
  \caption{Effect of calibration source on Wanda at 50\% unstructured
           sparsity (LLaMA-3-8B). The activation norm $\|X_j\|_2$ is a
           statistical fingerprint of the calibration distribution, so
           we expect domain-mismatched calibration to disproportionately
           damage capabilities not represented in that distribution.
           \textbf{[TO BE FILLED]}}
  \label{tab:exp1-calib-domain}
  \centering
  \small
  \begin{tabular}{lcccc}
    \toprule
    Method (calibration) & WT-2 ppl $\downarrow$ & MMLU $\uparrow$ & GSM8K $\uparrow$ & KoBEST HS $\uparrow$ \\
    \midrule
    Dense (no pruning)         & -- & -- & -- & -- \\
    \midrule
    Magnitude (no calibration) & -- & -- & -- & -- \\
    SparseGPT (C4)             & -- & -- & -- & -- \\
    \midrule
    Wanda (C4, English web)    & -- & -- & -- & -- \\
    Wanda (MC4-ko, Korean)     & -- & -- & -- & -- \\
    Wanda (The Stack, code)    & -- & -- & -- & -- \\
    Wanda (random tokens)      & -- & -- & -- & -- \\
    \bottomrule
  \end{tabular}
\end{table}

\subsection{Experiment 2: Per-Task Fragility Curves Across Sparsity}
\label{sec:exp2}

We hold calibration fixed (C4) and sweep sparsity at 50\%, 60\%, and 70\%.
The 50\% row for each method is shared with Table~\ref{tab:exp1-calib-domain},
so this table contributes six \emph{new} pruning runs.
The goal is to identify whether different downstream tasks cliff-drop at
different sparsity thresholds, and whether perplexity tracks those drops.

\begin{table}[h]
  \caption{Per-task fragility under increasing sparsity, LLaMA-3-8B with
           C4 calibration where applicable. \textbf{[TO BE FILLED]}}
  \label{tab:exp2-sparsity-curves}
  \centering
  \small
  \begin{tabular}{llcccc}
    \toprule
    Method & Sparsity & WT-2 ppl $\downarrow$ & MMLU $\uparrow$ & GSM8K $\uparrow$ & KoBEST HS $\uparrow$ \\
    \midrule
    Dense                & 0\%   & -- & -- & -- & -- \\
    \midrule
    Magnitude            & 50\%  & -- & -- & -- & -- \\
    Magnitude            & 60\%  & -- & -- & -- & -- \\
    Magnitude            & 70\%  & -- & -- & -- & -- \\
    \midrule
    SparseGPT (C4)       & 50\%  & -- & -- & -- & -- \\
    SparseGPT (C4)       & 60\%  & -- & -- & -- & -- \\
    SparseGPT (C4)       & 70\%  & -- & -- & -- & -- \\
    \midrule
    Wanda (C4)           & 50\%  & -- & -- & -- & -- \\
    Wanda (C4)           & 60\%  & -- & -- & -- & -- \\
    Wanda (C4)           & 70\%  & -- & -- & -- & -- \\
    \bottomrule
  \end{tabular}
\end{table}

\subsection{Experiment 3: Calibration $\times$ Sparsity Interaction}
\label{sec:exp3}

We test the extreme corner of the design space: English-calibrated and
Korean-calibrated Wanda, both at 50\% and 70\% sparsity. Three of the four
cells are reused from Tables~\ref{tab:exp1-calib-domain} and
\ref{tab:exp2-sparsity-curves}; only \emph{Wanda (MC4-ko) at 70\%} is a new
run. If the Korean-capability loss from English calibration and from high
sparsity compound super-additively on KoBEST HellaSwag, this would suggest
both effects target the same low-activation Korean features.

\begin{table}[h]
  \caption{Compound effect of calibration mismatch and high sparsity on
           Wanda (LLaMA-3-8B). $\Delta$ columns report drop relative to
           the dense baseline (filled after results are available).
           \textbf{[TO BE FILLED]}}
  \label{tab:exp3-calib-x-sparsity}
  \centering
  \small
  \begin{tabular}{llcccc}
    \toprule
    Calibration & Sparsity & MMLU $\uparrow$ & $\Delta$MMLU & KoBEST HS $\uparrow$ & $\Delta$KoBEST \\
    \midrule
    -- (Dense)         & 0\%   & -- & 0 & -- & 0 \\
    \midrule
    C4 (English)       & 50\%  & -- & -- & -- & -- \\
    C4 (English)       & 70\%  & -- & -- & -- & -- \\
    \midrule
    MC4-ko (Korean)    & 50\%  & -- & -- & -- & -- \\
    MC4-ko (Korean)    & 70\%  & -- & -- & -- & -- \\
    \bottomrule
  \end{tabular}
\end{table}

\subsection{Experiment 4 (Supplement): English-Only Domain Shift on LLaMA-2-7B}
\label{sec:exp4}

To isolate the cross-lingual component of Experiment 1, we run a parallel
domain-shift study on LLaMA-2-7B---which has limited Korean capability---using
only English-domain calibrations. If within-English domain shift (web text
vs.\ code vs.\ random tokens) produces \emph{less} damage than the
cross-lingual shift in Experiment 1, that supports the thesis that the
language axis dominates the domain axis for Wanda.

\begin{table}[h]
  \caption{English-only domain shift on LLaMA-2-7B. All evaluations are on
           English benchmarks; Magnitude and SparseGPT are reported only at
           the two corner sparsities for compute reasons.
           \textbf{[TO BE FILLED]}}
  \label{tab:exp4-en-only}
  \centering
  \small
  \begin{tabular}{llccc}
    \toprule
    Method (calibration)  & Sparsity & WT-2 ppl $\downarrow$ & MMLU $\uparrow$ & ARC-c $\uparrow$ \\
    \midrule
    Dense                 & 0\%  & -- & -- & -- \\
    \midrule
    Magnitude             & 50\% & -- & -- & -- \\
    Magnitude             & 70\% & -- & -- & -- \\
    SparseGPT (C4)        & 50\% & -- & -- & -- \\
    SparseGPT (C4)        & 70\% & -- & -- & -- \\
    \midrule
    Wanda (C4, web)       & 50\% & -- & -- & -- \\
    Wanda (C4, web)       & 60\% & -- & -- & -- \\
    Wanda (C4, web)       & 70\% & -- & -- & -- \\
    Wanda (The Stack)     & 50\% & -- & -- & -- \\
    Wanda (The Stack)     & 60\% & -- & -- & -- \\
    Wanda (The Stack)     & 70\% & -- & -- & -- \\
    Wanda (random tokens) & 50\% & -- & -- & -- \\
    Wanda (random tokens) & 60\% & -- & -- & -- \\
    Wanda (random tokens) & 70\% & -- & -- & -- \\
    \bottomrule
  \end{tabular}
\end{table}

\paragraph{Run accounting.}
Tables \ref{tab:exp1-calib-domain}--\ref{tab:exp4-en-only} together cover
$6 + 6 + 1 + 13 = 26$ distinct pruning runs; the dense rows do not count as
pruning runs. Cells reused across tables (e.g., Wanda(C4) at 50\% in both
Tables~\ref{tab:exp1-calib-domain} and \ref{tab:exp2-sparsity-curves}) are
counted once.

% =============================================================
\section{Analysis}
% =============================================================

\subsection{Calibration Domain Sensitivity}

\textbf{[To be written after experiments.]}
We hypothesize that English-calibrated Wanda will damage Korean capabilities
(KoBEST HellaSwag) more than Korean-calibrated Wanda damages English
capabilities (MMLU, GSM8K), because Korean text naturally contains some
English tokens but English text contains essentially no Korean. The decisive
prediction is that English-calibrated Wanda on KoBEST HellaSwag falls
\emph{at or below} the calibration-free magnitude pruning baseline---if so,
calibration data has been actively counterproductive.

\subsection{Sparsity Cliff and Per-Task Fragility}

\textbf{[To be written after experiments.]}
We expect WikiText-2 perplexity to degrade approximately continuously,
while downstream tasks cliff-drop at distinct thresholds---likely GSM8K and
KoBEST HellaSwag before MMLU. This would imply that perplexity is a poor
proxy for downstream capability retention under aggressive pruning.

\subsection{Calibration $\times$ Sparsity Interaction}

\textbf{[To be written after experiments.]}
If the Korean-capability loss from English calibration (Experiment 1) and
from high sparsity (Experiment 2) compound super-additively at 70\% sparsity
on KoBEST HellaSwag, this would suggest both effects damage the same set of
low-activation Korean-specific features, rather than two independent failure
modes.

\subsection{Cross-Lingual vs.\ Within-English Domain Shift}

\textbf{[To be written after experiments.]}
The LLaMA-2-7B supplement (Section~\ref{sec:exp4}) lets us compare the
within-English domain shift (web text $\rightarrow$ code $\rightarrow$ random
tokens) against the cross-lingual shift (English $\rightarrow$ Korean) from
Experiment 1. We expect the cross-lingual shift to produce strictly more
damage on the affected language's benchmarks, since dormant features are
fully missed rather than merely under-weighted.

% =============================================================
\section{Limitations}
% =============================================================

Our study has the following limitations:

\begin{itemize}
  \item \textbf{Scope}: Due to compute constraints, we do not evaluate
        LLaMA-30B or larger models.
  \item \textbf{Rigor}: Each experiment is run with a single random seed;
        variance across seeds is not reported.
  \item \textbf{Baselines}: We do not compare against more recent pruning
        methods (e.g., RIA~\cite{ria}) due to time constraints.
  \item \textbf{Hardware comparison}: Hardware-level benchmarks (latency,
        throughput) for actual sparse inference are not included, as
        unstructured sparsity does not directly translate to speedup on
        standard GPUs without specialized kernels.
\end{itemize}

% =============================================================
\section{Conclusion}
% =============================================================

We conducted a systematic stress-test of Wanda, a simple yet effective
LLM pruning method, along four axes: model architecture, calibration data,
sparsity ratio, and evaluation metrics. We additionally documented four
reproducibility issues in the original codebase arising from library version
updates. \textbf{[Conclusions to be updated after experiments.]}
Our results suggest that while Wanda is a strong baseline under its original
evaluation conditions, its performance is sensitive to \textbf{[X, Y, Z]},
which should be considered when applying the method in practice.

% =============================================================
% References
% =============================================================
\section*{References}
{
\small
[1] \label{wanda} Sun, M., Liu, Z., Bair, A., \& Kolter, J.Z. (2024).
A Simple and Effective Pruning Approach for Large Language Models.
\textit{ICLR 2024}. \url{https://arxiv.org/abs/2306.11695}

[2] \label{llama} Touvron, H., et al. (2023).
LLaMA: Open and Efficient Foundation Language Models.
\textit{arXiv:2302.13971}.

[3] \label{opt} Zhang, S., et al. (2022).
OPT: Open Pre-trained Transformer Language Models.
\textit{arXiv:2205.01068}.

[4] \label{sparsegpt} Frantar, E., \& Alistarh, D. (2023).
SparseGPT: Massive Language Models Can Be Accurately Pruned in One Shot.
\textit{ICML 2023}.

[5] \label{c4} Raffel, C., et al. (2020).
Exploring the Limits of Transfer Learning with a Unified Text-to-Text Transformer.
\textit{JMLR}.

[6] \label{ria} Zhang, M., et al. (2024).
Plug-and-Play: An Efficient Post-training Pruning Method for Large Language Models.
\textit{ICLR 2024}.
}

%%%%%%%%%%%%%%%%%%%%%%%%%%%%%%%%%%%%%%%%%%%%%%%%%%%%%%%%%%%%
\newpage
\section*{NeurIPS Paper Checklist}

\begin{enumerate}

\item {\bf Claims}
    \item[] Answer: \answerYes{}
    \item[] Justification: The abstract and introduction clearly state that
    this is a stress-testing study of Wanda, not a proposal of a new method.
    Claims are scoped accordingly.

\item {\bf Limitations}
    \item[] Answer: \answerYes{}
    \item[] Justification: Section 6 explicitly discusses limitations in scope,
    rigor, baselines, and hardware evaluation.

\item {\bf Theory assumptions and proofs}
    \item[] Answer: \answerNA{}
    \item[] Justification: This paper does not include theoretical results.

\item {\bf Experimental result reproducibility}
    \item[] Answer: \answerYes{}
    \item[] Justification: All commands, library versions, and hardware
    configurations are documented. Our modified code is publicly available.

\item {\bf Open access to data and code}
    \item[] Answer: \answerYes{}
    \item[] Justification: Our fork with all modifications is available at
    \url{https://github.com/[your-account]/wanda}.

\item {\bf Experimental setting/details}
    \item[] Answer: \answerYes{}
    \item[] Justification: Section 4 describes all models, datasets,
    hyperparameters, and evaluation metrics.

\item {\bf Experiment statistical significance}
    \item[] Answer: \answerNo{}
    \item[] Justification: Due to compute constraints, experiments are run
    with a single seed. Error bars are not reported.

\item {\bf Experiments compute resources}
    \item[] Answer: \answerYes{}
    \item[] Justification: Section 4.1 reports the GPU type, VRAM, CUDA
    version, and software stack used for all experiments.

\item {\bf Code of ethics}
    \item[] Answer: \answerYes{}
    \item[] Justification: This work studies model compression of
    publicly available models and poses no ethical concerns.

\item {\bf Broader impacts}
    \item[] Answer: \answerNA{}
    \item[] Justification: This is a reproducibility study with no direct
    negative societal impact.

\item {\bf Safeguards}
    \item[] Answer: \answerNA{}
    \item[] Justification: We do not release new models or datasets.

\item {\bf Licenses for existing assets}
    \item[] Answer: \answerYes{}
    \item[] Justification: Wanda (MIT License), OPT (OPT-175B License),
    LLaMA-2 (LLaMA 2 Community License), C4 (ODC-BY) are all properly cited.

\item {\bf New assets}
    \item[] Answer: \answerNA{}
    \item[] Justification: No new datasets or models are introduced.

\item {\bf Crowdsourcing and research with human subjects}
    \item[] Answer: \answerNA{}
    \item[] Justification: No human subjects are involved.

\item {\bf Institutional review board (IRB) approvals}
    \item[] Answer: \answerNA{}
    \item[] Justification: No human subjects are involved.

\item {\bf Declaration of LLM usage}
    \item[] Answer: \answerNA{}
    \item[] Justification: LLMs are the \emph{subject} of study, not
    a component of the methodology.

\end{enumerate}

\end{document}